\section{Introdução}

\subsection{Contextualização}

O avanço recente dos grandes modelos de linguagem (large language models -- LLMs) ampliou significativamente o interesse por soluções capazes de responder perguntas em linguagem natural \parencite{minaee2024llmSurvey}. No setor público, esse potencial mostra-se particularmente relevante em contextos caracterizados por acervos extensos, dinâmicos e heterogêneos, tais como legislação, pareceres, notas técnicas, discursos parlamentares e documentos administrativos \parencite{ipu2024aiParliaments}. Nesses ambientes, o desafio não se restringe ao armazenamento da informação, mas abrange, sobretudo, sua recuperação qualificada, contextualizada e verificável \parencite{wfd2023guidelinesAiParliaments}.

Apesar de seu potencial, o uso direto de LLMs apresenta limitações amplamente discutidas na literatura, entre as quais se destacam a alucinação, a desatualização do conhecimento usado no treinamento e a dificuldade de identificar as evidências utilizadas na produção das respostas \parencite{gao2023ragSurvey}. Tais limitações assumem maior gravidade em contextos institucionais nos quais se exige transparência, verificabilidade, controle das fontes e responsabilização sobre os conteúdos produzidos \parencite{wfd2023guidelinesAiParliaments}.

Nesse cenário, a abordagem retrieval-augmented generation (RAG) apresenta-se como alternativa tecnicamente promissora, pois, em vez de depender exclusivamente do conhecimento internalizado no modelo, sistemas RAG recuperam trechos relevantes de uma base documental externa e os incorporam ao contexto da geração \parencite{gao2023ragSurvey}. Com isso, tende-se a elevar a precisão factual e a atualidade das respostas, ao mesmo tempo em que se favorece a explicitação das fontes que as sustentam \parencite{gao2023ragSurvey}.

No contexto do Parlamento brasileiro, a organização do conhecimento legislativo e a ampliação do acesso qualificado à informação constituem desafios relevantes de pesquisa aplicada, como observa \textcite{martelloteViola2025organizacaoConhecimento} ao defender que o emprego de inteligência artificial nesse domínio deve ser acompanhado de atenção à estrutura informacional, à mediação documental e à explicabilidade dos resultados.

É nesse horizonte que se insere o presente trabalho, ao propor o desenvolvimento de um chatbot legislativo baseado em RAG, voltado à consulta de discursos da 56\textordfeminine{} Legislatura do Senado Federal.

\subsection{Motivação}

A administração pública contemporânea opera sob crescente pressão por transparência, rastreabilidade, eficiência e capacidade de resposta \parencite{un2015responsiveGovernance}. No ambiente parlamentar, a consulta a discursos e pronunciamentos mostra-se relevante para diversas finalidades, tais como análise temática, comparação de posicionamentos, apoio técnico a gabinetes, elaboração de relatórios, produção de subsídios e atendimento ao cidadão \parencite{skubicFiser2024discourseReview}. Entretanto, a simples disponibilização pública desses documentos não garante, por si só, acesso eficiente ao conteúdo neles contido \parencite{geunis2023parliamentaryMonitoring}.

Em grandes acervos discursivos, mecanismos tradicionais de busca, predominantemente baseados em correspondência lexical, tendem a ser insuficientes quando o usuário formula perguntas complexas, comparativas, interpretativas ou orientadas à síntese \parencite{gupta2024comprehensiveRag}. Além disso, a recuperação da informação frequentemente depende de conhecimento prévio sobre vocabulário, datas, autores ou a estrutura da base documental. Como consequência, a informação permanece formalmente acessível, mas substantivamente difícil de explorar.

Paralelamente, soluções baseadas exclusivamente em LLMs generalistas podem produzir respostas plausíveis e linguisticamente bem estruturadas, porém desprovidas de fundamentação verificável \parencite{anhHoang2025hallucinations,openai2025hallucinate}. Em domínios institucionais, tal comportamento é particularmente problemático, pois compromete a confiança, a responsabilização e o uso controlado da tecnologia \parencite{highberg2025responsibleAi}. \textcite{deene2025raggdpr} argumenta que sistemas RAG podem mitigar parte desses riscos ao deslocar a base do conhecimento para documentos controlados e atualizáveis. Na mesma direção, \textcite{said2025ragPractice} enfatiza que, em cenários de produção, sistemas RAG requerem práticas adicionais de versionamento, observabilidade e avaliação contínua, de modo a assegurar qualidade e consistência ao longo do tempo.

Dessa forma, a motivação deste trabalho decorre da necessidade de investigar uma solução tecnicamente viável e institucionalmente adequada para ampliar o acesso qualificado a discursos parlamentares, combinando busca semântica, geração textual e verificação documental.

\subsection{Problema de pesquisa}

Diante desse contexto, formula-se o seguinte problema de pesquisa: como estruturar uma solução baseada em RAG capaz de permitir a consulta, em linguagem natural, a discursos da 56\textordfeminine{} Legislatura do Senado Federal, produzindo respostas úteis, fundamentadas em evidências recuperadas e compatíveis com exigências de verificação em contexto institucional?

A formulação desse problema decorre da percepção de que o acesso documental, embora formalmente assegurado, ainda enfrenta limitações relevantes quando se exige recuperação semântica, síntese orientada ao usuário e indicação clara das fontes utilizadas na resposta.

\subsection{Objetivos}

O objetivo geral deste trabalho consiste em desenvolver e demonstrar uma solução de chatbot legislativo baseada em RAG para consulta a discursos parlamentares da 56\textordfeminine{} Legislatura do Senado Federal.

Como objetivos específicos, pretende-se:
\begin{enumerate}
  \item estruturar uma base de conhecimento consultável a partir dos discursos parlamentares;
  \item possibilitar a formulação de perguntas em linguagem natural com respostas fundamentadas em trechos recuperados;
  \item avaliar empiricamente a qualidade inicial dos processos de recuperação e geração;
  \item documentar uma arquitetura reprodutível, modular e passível de evolução para cenários institucionais.
\end{enumerate}

\subsection{Justificativa}

A relevância deste estudo decorre de sua contribuição em dois planos complementares. No plano prático, a pesquisa investiga uma forma de ampliar o acesso qualificado a informações parlamentares públicas, contribuindo para transparência, memória institucional e apoio a atividades de gabinete, pesquisa e controle social. No plano teórico-aplicado, o trabalho examina a aplicação de técnicas contemporâneas de recuperação e geração em um domínio documental complexo, marcado por exigências elevadas de explicabilidade, confiabilidade e vinculação às fontes.

Além disso, o estudo se justifica pela aderência do problema a desafios concretos da administração pública digital. Em vez de tratar a inteligência artificial como mecanismo autônomo de produção de respostas, propõe-se abordagem em que a geração textual é condicionada por fontes documentais públicas, controláveis e auditáveis. Essa orientação é coerente com demandas institucionais por uso responsável de IA em contextos sensíveis.
